\subsection{Hard Negative Mining}

Following the initial training phase of the bi-encoder, a hard negative mining procedure was utilized to generate additional training examples for continued fine-tuning. The procedure was introduced to provide more challenging negative examples than those obtained through random sampling.

The trained bi-encoder was used to encode the regulatory entity collection into dense vector representations, which were indexed for similarity-based retrieval. For each regulatory reference, the corresponding embedding was generated and used to retrieve the top-$k$ most similar regulatory entities.

The ground-truth regulatory entity was removed from the retrieved candidate set, and the highest-ranked remaining entities were selected as hard negatives. These candidates represented semantically similar but incorrect regulatory entities, such as documents with overlapping titles, related regulatory terminology, or similar identifiers.

To prevent incorrect negative assignments, retrieved candidates were compared against the labelled reference-entity associations generated during dataset construction. Candidates corresponding to valid reference-entity relationships were excluded before being added to the hard negative set.

The resulting hard negative examples were incorporated into the training dataset and used during continued fine-tuning of the bi-encoder. This expanded training procedure allowed the model to learn improved separation between closely related regulatory entities.